\clearpage\fancyhead[L]{\small BORDA’S KREIS.}\hypertarget{native-409-part-lower}{}\label{native-409-part-lower}
\begin{center}\large BORDA’S KREIS.\end{center}
\begin{center}\hrule\vspace{4pt}\fontsize{10}{12}\selectfont V. \textsc{Liadenau} u. \textsc{Bohnenberger}, Zeitschrift für Astronomie u. verwandte Wiss.\\ Band V. Seite 198..211. März und April 1818.\par\vspace{4pt}\hrule\vspace{4pt}\end{center}
\par\indent Wie von Ihnen gewünscht, hochverehrter Herr Professor, freue ich mich, Ihnen hier die Methode mitzuteilen, mit der ich die Visierlinie parallel zur Kreisebene eingestellt habe. Wie mich die wiederholte Erfahrung gelehrt hat, lässt sich damit eine sehr hohe Genauigkeit erzielen.

\par\indent Um mich kurz und genau ausdrücken zu können, bezeichne ich die verschiedenen Achsen des Instruments wie folgt. Die Achse, die bei Höhenmessungen vertikal gestellt werden muss und die man auch bei jeder anderen Verwendung zumindest annähernd in dieser Lage belässt, nenne ich die \emph{erste Achse}. Die Bewegung um diese Achse wird am Azimutkreis gemessen. Dieser hat jedoch nur einen Index, dessen Nonius nur Minuten angibt, obwohl sich noch Bruchteile einer Minute schätzen lassen. Außerdem ist ungewiss, ob eine kleine Exzentrizität vorliegt. Eine Berichtigung der Visierlinie mithilfe dieses Azimutkreises (indem man in zwei entgegengesetzten vertikalen Lagen des Hauptkreises das Azimut eines irdischen Gegenstands bestimmt, nötigenfalls unter Berücksichtigung seiner Parallaxe, wenn er nicht sehr weit entfernt ist) liefert daher noch nicht die ganze erreichbare oder wünschenswerte Genauigkeit. Für eine erste vorläufige Berichtigung kann man sich jedoch dieses bequemen und einfachen Verfahrens bedienen.

\par\indent Fast senkrecht zur ersten Achse liegt diejenige, die ich die \emph{zweite} nenne. Sie liegt also stets nahezu horizontal. Durch Drehung um sie wird der Kreis in eine vertikale, geneigte oder horizontale Lage gebracht. An den \textsc{Reichenbach}schen Kreisen gibt es keine Teilung, mit der sich das Ausmaß dieser Drehung messen lässt. Für die meisten Beobachtungen ist sie auch überflüssig, aber es gibt doch einige Fälle, in denen sie wünschenswert wäre. Insbesondere könnte man bei Tage mühelos Winkel zwischen Fixsternen und irdischen Objekten messen, was gegenwärtig entweder gar nicht oder nur mit großer Schwierigkeit möglich ist.

\par\indent Wiederum fast senkrecht zur zweiten Achse liegt die Achse, um die sich der ganze Kreis dreht. Diese Bewegung wird am Trommelkreis gemessen. Da diese Achse im mathematischen Sinn mit derjenigen zusammenfällt, \SourceLeaf{410} um die sich der Alhidadenkreis bewegt, werden beide Achsen als eine betrachtet, die ich die \emph{dritte} nenne. Diese dritte Achse steht genau senkrecht auf der Kreisebene. Bezeichnet $a$ den Winkel der Visierlinie mit der Kreisebene (den ich positiv ansetze, wenn die Objektivseite weiter von der Kreisfläche entfernt ist als die Okularseite), so ist $90^\circ{-}a$ der Winkel der Visierlinie mit der dritten Achse, wobei jene zur Objektivseite hin und diese von der graduierten Fläche aus abwärts genommen wird.

\par\indent Für die Benutzung des Instruments braucht die zweite Achse weder genau parallel zur Kreisebene noch genau senkrecht zur ersten Achse zu stehen. Die Abweichungen werden jedoch gewiss klein sein. Wiederholte Erfahrungen mit dem hiesigen 12-zölligen \textsc{Reichenbach}schen Kreis haben mich gelehrt, dass sie jahrelang unverändert bleiben, selbst nachdem das Instrument auseinandergenommen worden ist. Mit $b$ bezeichne ich die Neigung der zweiten Achse gegen die Kreisebene, positiv genommen, wenn ihr linker Arm\footnote{wenn man hinter dem als vertikal vorausgesetzten Kreis steht.} näher am Kreis liegt. Schließlich bezeichne ich mit $c$ die Neigung der zweiten Achse gegen eine Ebene senkrecht zur ersten Achse (als solche kann man die Ebene des Azimutkreises ansehen), positiv genommen, wenn der linke Arm höher über ihr liegt.

\par\indent Ferner seien $M$ und $N$ die Schnittpunkte der dritten Achse mit der zweiten Achse und mit der Visierlinie und $O$ ein gut sichtbares Objekt. Der Abstand $MN$ hat bei jedem Kreis einen bestimmten, durch Messung oder auf andere Weise zu ermittelnden Wert. Beim hiesigen Kreis beträgt er $\frac{1}{8}$ Meter. Ist die Visierlinie auf das Objekt $O$ gerichtet, bezeichne ich den Winkel $MON$ mit $p$. Die genaue Formel dafür ist

\[ \sin p = \frac{MN}{MO}\cdot\cos a \]

\noindent wofür man jederzeit $p = \dfrac{MN}{MO}\cdot 206265''$ ansetzen darf. Es wird nicht schwer sein, $MO$ so genau zu bestimmen, dass $p$ auf die Sekunde genau bekannt ist. Bei weit entfernten Objekten wird $p$ ganz unmerklich. Das ist aber kein Grund, sie bei der folgenden Berichtigungsmethode näheren Objekten vorzuziehen, wenn letztere bessere Zielpunkte bieten.

\par\indent Ich setze voraus, dass bei allen hier zu beschreibenden Arbeitsschritten der Kreis selbst überhaupt nicht um die dritte Achse bewegt wird, dass also die Klemmung an der Trommel fest angezogen ist und ihre Stellschraube unberührt bleibt. Man kann, wenn man will, auch die Pressmutter am hinteren Teil der dritten Achse etwas fester anziehen als gewöhnlich.

\par\indent Der Kürze halber bezeichne ich den Winkel zwischen der Ebene durch die zweite und dritte Achse und der Ebene durch die erste und zweite Achse mit $I$. Er bestimmt somit das Ausmaß der Drehung des Kreises um die zweite Achse. Die Bauart des Instruments erlaubt nur Werte dieses Winkels von $0$ bis $90^\circ$ und etwas über beide Grenzen hinaus. Dass die Reichenbachschen Kreise keine Teilung zur Messung dieses Winkels haben, habe ich schon oben bemerkt.

\par\indent Das Wesentliche meiner Methode liegt darin, dass sich die Visierlinie, wenn sie bei irgendeiner Lage des Instruments auf ein Objekt gerichtet ist, durch die Verbindung einer Drehung um die zweite Achse mit einer Drehung des Alhidadenkreises um die dritte ein zweites Mal auf das Objekt richten lässt. Ich nehme an, dass $I = i$ ist, wenn die dritte Achse mit der zweiten Achse und der Geraden $MO$ in eine Ebene gebracht ist. Dagegen soll der Index des Alhidadenkreises $k$ anzeigen, wenn die Visierlinie mit der zweiten und dritten Achse in einer Ebene liegt, sodass die Visierlinie und die zweite Achse (linker Arm) auf der Kreisebene dieselbe Projektion haben

\SourceLeaf{411}
\par\indent Man sieht leicht, dass, wenn bei auf das Objekt gerichteter Visierlinie der Index des Alhidadenkreises auf $k-\varkappa$ steht und $I=i-\eta$ ist, die Visierlinie wieder auf dasselbe Objekt gerichtet ist, wenn man $I=i+\eta$ einstellt und den Index des Alhidadenkreises auf $k+\varkappa$ stellt. Diese zweite Stellung ist natürlich nur dann physisch möglich, wenn der zweite Wert von $I$ innerhalb der Grenzen liegt, die die Bauart des Instruments zulässt. Ferner lässt sich leicht zeigen, dass die folgende Gleichung exakt gilt:

\[ \sin b \cos\varkappa = \cos b\, \mathrm{tang}(a+p) + \mathrm{cotang}\,\eta \sin\varkappa \]

\noindent und dass der Kosinus des Winkels zwischen der zweiten Achse und der Geraden $MO$

\[ = \sin b \sin(a+p) + \cos b \cos(a+p) \cos\varkappa \]

\noindent wird. Daraus folgt leicht, dass $\varkappa$ entweder sehr klein ist oder sehr nahe bei $180^\circ$ liegt. Im ersten Fall ist der genannte Winkel sehr klein, im zweiten liegt er nahe bei $180^\circ$. Das heißt, das Objekt muss im ersten Fall annähernd in Richtung des linken Arms der zweiten Achse liegen, im zweiten annähernd in Richtung des rechten.

\par\indent Damit erhält man hinreichend genau

\noindent im ersten Fall \[ b - a - p = \frac{\varkappa}{\mathrm{tang}\,\eta} \]

\noindent im zweiten Fall \[ b + a + p = \frac{\varkappa - 180^\circ}{\mathrm{tang}\,\eta} \]

\par\indent Aus zwei solchen Versuchen erhält man also, sofern man $\eta$ ermitteln kann, die Werte von $b-a-p$ und von $b+a+p$ und folglich von $b$ und $a+p$ oder, sofern $p$ als bekannt gelten kann, von $a$. Zugleich erhält man offenbar aus jedem solchen Versuch, ohne $\eta$ zu benötigen, eine Bestimmung von $k$. Da $\eta$ hierbei frei gewählt werden kann, macht man es am besten so groß wie möglich. Allerdings wird man kaum weit über $45^\circ$ hinausgehen können. Da man an \textsc{Reichenbach}schen Kreisen die Werte von $I$ nicht ablesen kann, muss man sich auf solche Werte beschränken, für die man die Stellung auf anderem Wege erreichen kann. Ich werde weiter unten zeigen, wie man die Stellung, für die $I=0$ ist, in der also die erste. zweite und dritte Achse in einer Ebene liegen, mit größter Genauigkeit einstellen kann. Ferner kann man die Kreisebene zur ersten Achse parallel machen, indem man beide auf bekannte Weise vertikal stellt. Diese beiden Stellungen kann man durch zwei Striche auf dem Quadranten markieren, der die Klemmung für die Drehung um die zweite Achse trägt. Die Striche zieht man an der Kante einer der Stützen der zweiten Achse, damit man die Stellungen jederzeit schnell und hinreichend genau wiederfindet. Der Kürze halber nenne ich die erste Stellung die \emph{horizontale}, die zweite die \emph{vertikale}. Für die letztere gilt exakt

\[ \cos I = -\,\mathrm{tang}\,b\,.\,\mathrm{tang}\,c \]

\noindent $I$ wird also nicht merklich von $90^\circ$ abweichen, und da der Unterschied $=2\eta$ gesetzt ist, ergibt sich exakt

\[ \mathrm{tang}\,\eta = \BareRoot{\dfrac{1+\mathrm{tang}\,b\,.\,\mathrm{tang}\,c}{1-\mathrm{tang}\,b\,.\,\mathrm{tang}\,c}} = \BareRoot{\dfrac{\cos(b-c)}{\cos(b+c)}} \]

\noindent und damit wird der Wert nur unmerklich von $\eta$ abweichen.

\SourceLeaf{412}
\par\indent Daraus folgt: Wenn beim ersten Versuch, bei dem das Objekt annähernd in Richtung des linken Arms der zweiten Achse lag, der Index des Alhidadenkreises in der horizontalen Stellung des Kreises $h$ und in der vertikalen $\upsilon$ angezeigt hat, und wenn er dagegen beim zweiten Versuch, bei dem das Objekt annähernd in Richtung des rechten Arms der zweiten Achse lag, entsprechend $180^\circ+h'$, $180^\circ+\upsilon'$ angezeigt hat, erhält man:

\[ \begin{aligned} b-a-p &= \upsilon-k = k-h = \tfrac{1}{2}(\upsilon-h) \\ b+a+p &= \upsilon'-k = k-h' = \tfrac{1}{2}(\upsilon'-h') \end{aligned} \]

\noindent also

\[ \begin{aligned} a &= \tfrac{1}{4}(\upsilon'-h'-\upsilon+h)-p \\ b &= \tfrac{1}{4}(\upsilon'-h'+\upsilon-h) \end{aligned} \]

\noindent Zugleich erhält man für $k$ die beiden Werte

\[ \tfrac{1}{2}(\upsilon+h) \quad\text{und}\quad \tfrac{1}{2}(\upsilon'+h') \]

\noindent deren enge Übereinstimmung die Richtigkeit der Arbeitsschritte bestätigen wird.

\par\indent Mein Verfahren läuft demnach wie folgt ab. Das gewählte Objekt muss einen scharf erkennbaren Zielpunkt bieten und darf höchstens einige Grade vom Horizont entfernt sein. Ich stelle das Instrument so auf, dass einer seiner Füße annähernd in einer Ebene mit der ersten Achse und dem Objekt liegt. In diese Ebene bringe ich dann auch die zweite Achse durch eine Drehung im Azimut nach Augenmaß, sodass ihr linker Arm zunächst dem Objekt zugewandt ist. Danach stelle ich den Kreis horizontal und richte die Visierlinie mit der Schraube an dem genannten Fuß und durch Verstellen der Alhidade auf das Objekt. Dann wird der Kreis vertikal gestellt und die Visierlinie durch Verstellen im Azimut und an der Alhidade wieder auf das Objekt gerichtet. Damit ist die erste grobe Einstellung abgeschlossen. Der Kreis wird erneut horizontal gestellt. Wenn das Verstellen der Alhidade allein nicht genügt, um die Visierlinie auf das Objekt zu richten, hilft man an der Fußschraube nach. Dann wird der Kreis wieder vertikal gestellt und die Visierlinie durch Verstellen der Alhidade auf das Objekt gerichtet, nötigenfalls mit Hilfe der Azimutschraube. Man fährt mit diesen beiden Arbeitsschritten abwechselnd fort, bis das Verstellen der Alhidade allein in horizontaler Lage zum genauen Anvisieren genügt. Die Ablesungen am Alhidadenkreis ergeben $h$ und $\upsilon$.

\par\indent Danach löst man nach der letzten vertikalen Einstellung die Azimutklemmung und dreht den Kreis um die erste Achse um $180^\circ$, sodass nun der rechte Arm der zweiten Achse annähernd auf das Objekt gerichtet ist. Durch Verstellen der Alhidade (um nahezu $180^\circ$) und eine Feineinstellung an der Azimutschraube richtet man die Visierlinie genau auf das Objekt. Auch hier lese ich den Index des Alhidadenkreises ab, der etwa $180^\circ+V$ anzeigen möge. Damit beginnt die zweite Reihe von Arbeitsschritten. Man legt den Kreis horizontal und visiert das Objekt durch Verstellen der Fußschraube und der Alhidade an. In der neuen vertikalen Stellung verstellt man im Azimut und an der Alhidade. So fährt man wieder abwechselnd fort, bis das Verstellen der Alhidade allein ausreicht, um das Objekt in beiden Lagen zu treffen. Die Ablesungen am Alhidadenkreis ergeben hier $180^\circ+h'$ und $180^\circ+\upsilon'$. Damit sind die Messungen abgeschlossen.

\par\indent Die zuvor genannte Ablesung $180^\circ+V$ dient auch dazu, die Größe $c$ zu bestimmen. Bei vertikaler Stellung des Kreises liegt die Visierlinie offenbar in einer Ebene mit der ersten und dritten Achse und ist nach oben gerichtet, wenn der Index des Alhidadenkreises auf $\tfrac{1}{2}(\upsilon+V)-90^\circ$ steht. Anders gesagt ist dann die Projektion der Visierlinie auf die Kreisebene parallel zur ersten Achse. Die Projektion der zweiten Achse (linker Arm) auf dieselbe Ebene bildet also mit dem zur ersten Achse parallelen Radius den Winkel $k+90^\circ-\tfrac{1}{2}(\upsilon+V)$. Daher gilt

\SourceLeaf{413}
\[ \sin c = \cos b\,\sin\bigl(\tfrac{1}{2}(\upsilon+V)-k\bigr) \]

\noindent wofür man

\[ c = \tfrac{1}{2}(\upsilon+V)-k = \tfrac{1}{2}(V-h) \]

\noindent setzen darf.

\par\indent Mit diesem Verfahren erhält man also den Fehler $a$ in Sekunden. Will man ihn bei positivem $a$ beseitigen, so muss das Fadenkreuz vom Kreis aus abwärts verschoben werden oder, wenn man das prismatische Okular verwendet, in Richtung vom Objektiv zum Okular. Um die Größe der Verschiebung unmittelbar vor Augen zu haben, kann man entweder das Verhältnis zu einem geeigneten Objekt von bekanntem Durchmesser nach Augenmaß schätzen oder zuerst die Visierlinie genau auf ein feines Ziel richten und dann den Alhidadenkreis um $a$ verschieben. Darauf verschiebt man das Fadenkreuz so weit, bis das Objekt in der Diagonale erscheint. Bei diesem heiklen Arbeitsschritt muss man selbstverständlich vorsichtig vorgehen und jede Verschiebung des Kreises selbst vermeiden. Durch Wiederholung des vorstehenden Verfahrens kann man sich stets vergewissern, ob man sein Ziel erreicht hat.

\par\indent Mir bleibt nun noch zu zeigen, wie man den Kreis horizontal einstellt, sodass $I=0$ ist. Solange $c$ und $b$ nicht beträchtlich sind, kommt es freilich nicht so genau darauf an, und man könnte sich mit dem Augenmaß begnügen. Mit dem folgenden Verfahren lässt sich die Stellung jedoch mit äußerster Genauigkeit erreichen.

\par\indent Nachdem bei feststehender Trommel $k$ bestimmt worden ist, legt man den Kreis zuerst nach Augenmaß horizontal, stellt den Index des Alhidadenkreises auf $k+l$, wobei $l$ einen beliebigen Winkel bezeichnet, am besten $90^\circ$, und bringt durch Verstellen im Azimut und am Fuß ein gut sichtbares Objekt auf die Visierlinie. Darauf stellt man den Index des Alhidadenkreises auf $k-l$. Hat man die horizontale Stellung zufällig genau getroffen, so lässt sich die Visierlinie durch Verstellen im Azimut allein wieder genau auf das Objekt richten. Andernfalls geht sie darüber oder darunter hinweg. Nachdem man das Objekt durch Verstellen im Azimut so nahe wie möglich an die Visierlinie gebracht hat (das heißt auf den zur dritten Achse parallelen Faden oder beim prismatischen Okular auf den zum Fernrohr parallelen Faden), bringt man es durch Drehen des Kreises um die zweite Achse um die Hälfte näher an die Visierlinie. Die andere Hälfte korrigiert man mit einer Fußschraube, nötigenfalls mit Hilfe der Azimutschraube, sodass das Objekt wieder genau auf der Visierlinie liegt. Um zu prüfen, wie genau halbiert wurde, stellt man den Index des Alhidadenkreises wieder auf $k+l$ und prüft, ob das Verstellen im Azimut allein das Objekt jetzt wieder genau auf die Visierlinie bringt. Falls nicht, korrigiert man wieder die Hälfte des jetzt gewiss viel kleineren Fehlers an der Schneckenschraube, die zur Feinbewegung um die zweite Achse dient. Man fährt damit fort, bis man ganz zufrieden ist.

\par\indent Ist die erste Einstellung nach Augenmaß zu schlecht ausgefallen, so ist die erste Halbierung schwer mit hinreichender Genauigkeit auszuführen, weil das Verstellen im Azimut das Objekt bei groß gewähltem $l$ gar nicht erst ins Gesichtsfeld bringt. Um sich also vieles Ausprobieren zu ersparen, wählt man $l$ anfangs lieber klein, etwa $=10^\circ$, und bringt den Kreis nach dem beschriebenen Verfahren in die gewünschte Stellung. Diese lässt sich dann durch Wiederholung des Verfahrens mit einem großen Wert von $l$ fein berichtigen.

\par\indent Wenn bei einem Instrument zufällig $b=c$ ist, steht die Kreisebene in der horizontalen Stellung genau senkrecht zur ersten Achse. Ein Objekt, das einmal auf der Visierlinie liegt, lässt sich nach jeder Änderung der Alhidadenstellung durch bloßes Verstellen im Azimut wieder auf sie bringen. \SourceLeaf{414} Sind aber $b$ und $c$ ungleich, so sieht man leicht, dass dies unmöglich ist. Die Kreisebene kann dann im \emph{eigentlichen} Sinne nur dann genau horizontal gestellt werden, wenn die erste Achse etwas geneigt wird. Beim hiesigen \textsc{Reichenbach}schen Kreis ist $c=+5'40''$, $b=+28''$ im Mittel aus mehreren Versuchen.

\par\indent Wie schon oben bemerkt, ist die Bestimmung von $k$ von der Kenntnis bestimmter Werte von $I$ ganz unabhängig. Ohne die Azimut- und Trommelschraube zu berühren, braucht man nur ein Objekt durch Verstellen der Alhidade und Drehen um die zweite Achse zweimal auf die Visierlinie zu bringen und den Mittelwert aus beiden Stellungen des Index des Alhidadenkreises zu nehmen. Am zweckmäßigsten richtet man es so ein, dass die Drehung um die zweite Achse möglichst groß wird. Dazu nehme ich zunächst die oben beschriebene grobe Einstellung vor. Wenn das Objekt dann bei ungefähr horizontaler Stellung des Kreises genau auf der Visierlinie liegt, drehe ich den Kreis um die zweite Achse, bis die Bogenbewegung des Objekts im Gesichtsfeld es wieder auf den zur Kreisebene parallelen Faden führt (beim prismatischen Okular auf den entsprechenden Faden, der also senkrecht zum Fernrohr steht). Dann bringe ich das Objekt durch Verstellen der Alhidade genau auf das Kreuz. \emph{Bliebe} das Objekt bei der Drehung um die zweite Achse genau auf dem Kreuz, so wäre dies ein Beweis dafür, dass es in der zweiten Achse selbst läge, dass folglich $b=a+p$ wäre und die Stellung des Index selbst $=k$.

\par\indent Bei dieser Gelegenheit möchte ich noch das Verfahren erwähnen, das ich mit Vorteil verwende, um die Kreisebene genau vertikal zu stellen. Dafür ist bei \textsc{Reichenbach}s Kreisen eigentlich die Hängelibelle vorgesehen, mit der sich zwar eine große Genauigkeit erzielen lässt. Ihre Verwendung ist aber etwas umständlich, weil das Fernrohr abgenommen werden muss, und erfordert eine sehr stabile Aufstellung und viel Vorsicht. Ich verwende den Quecksilberhorizont und wähle ein möglichst hoch gelegenes Objekt (es braucht übrigens nicht weit entfernt zu sein, denn bei diesem Arbeitsschritt schadet es nicht, wenn man das Okularrohr weit herausziehen muss, um es deutlich zu sehen). Ich stelle den Kreis nach Augenmaß vertikal, bringe ihn in die Vertikalebene des Objekts und visiere es genau an. Darauf wird das Fernrohr abgesenkt, bis das Bild des Objekts im Quecksilberhorizont zu sehen ist. Trifft die Visierlinie genau auf das Bild, so steht die Kreisebene vertikal. Andernfalls korrigiert man die eine Hälfte entweder an einer Fußschraube oder an der Schneckenschraube, die andere an der Azimutschraube. Hat man genau halbiert, trifft die Visierlinie des wieder angehobenen Fernrohrs das Objekt selbst genau. Andernfalls korrigiert man erneut die eine Hälfte der Abweichung an der Fußschraube oder der Schneckenschraube, die andere an der Azimutschraube. Ist die Höhe des Objekts nicht zu gering und lässt es sich genau genug anvisieren, so kann man auch auf diese Weise eine sehr große Genauigkeit erreichen, weit mehr als nötig ist, wenn man keine \emph{sehr} kleinen Zenitdistanzen messen will. Für \emph{allzu kleine} Zenitdistanzen empfiehlt es sich aber bekanntlich nicht, den \textsc{Borda}schen Kreis zu verwenden.

